\documentclass[
	fontsize=12pt,
	paper=a4,
	numbers=noenddot,
	parskip=half,
	openany
]{scrbook}
\usepackage{tcolorbox}

\include{../../../for_latex/preamble}
% ##################################################
% Gemeinsames Setup fuer die Probeklausuren
% (Java-Listings, Hinweis-/Beispiel-Boxen, Struktogramm-Stile in TikZ)
% wird nach \include{../for_latex/preamble} per \input eingebunden
% ##################################################

% --- Java Code-Highlighting ---
\definecolor{codegray}{rgb}{0.5,0.5,0.5}
\definecolor{codepurple}{rgb}{0.58,0,0.82}
\definecolor{backcolour}{rgb}{0.95,0.95,0.92}
\definecolor{keywordcolor}{rgb}{0.8,0.47,0.13}

\lstdefinestyle{javastyle}{
	backgroundcolor=\color{backcolour},
	commentstyle=\color{codegray},
	keywordstyle=\color{keywordcolor}\bfseries,
	numberstyle=\tiny\color{codegray},
	stringstyle=\color{codepurple},
	basicstyle=\ttfamily\footnotesize,
	breakatwhitespace=false,
	breaklines=true,
	captionpos=b,
	keepspaces=true,
	numbers=left,
	numbersep=5pt,
	showspaces=false,
	showstringspaces=false,
	showtabs=false,
	tabsize=2,
	language=Java,
	frame=single,
	rulecolor=\color{black!30}
}

\lstset{style=javastyle}

% --- Hinweis-/Beispiel-/Formel-Umgebungen (schlicht, ohne farbige Box) ---
\newenvironment{hinweis}[1][Hinweis]%
	{\par\smallskip\noindent\textbf{#1.}\enspace\ignorespaces}%
	{\par\smallskip}

\newenvironment{beispiel}[1][Beispiel]%
	{\par\smallskip\noindent\textbf{#1:}\par\nobreak\vspace{2pt}}%
	{\par\smallskip}

\newenvironment{formeln}[1][Gegebene Formeln]%
	{\par\smallskip\noindent\textbf{#1:}\par\nobreak\vspace{2pt}}%
	{\par\smallskip}

% ##################################################
% Struktogramm (Nassi-Shneiderman) in reinem TikZ
% ##################################################
\newlength{\nswidth}\setlength{\nswidth}{11.5cm}   % Gesamtbreite
\newlength{\nsind}\setlength{\nsind}{0.55cm}        % Einrueckung je Schleifenebene
% vorberechnete Versatzbreiten fuer Verzweigungen (in Koordinaten nutzbar)
\newlength{\nshalf}\setlength{\nshalf}{\dimexpr\nswidth/2\relax}            % halbe Gesamtbreite
\newlength{\nshalfi}\setlength{\nshalfi}{\dimexpr(\nswidth-\nsind)/2\relax}  % halbe Breite, 1x eingerueckt
\newlength{\nshalfii}\setlength{\nshalfii}{\dimexpr(\nswidth-2\nsind)/2\relax} % halbe Breite, 2x eingerueckt

\tikzset{
	% Basis fuer Anweisungs-/Verzweigungsbloecke
	nb/.style={draw, line width=0.4pt, anchor=north west, inner sep=4pt,
		align=left, font=\footnotesize, fill=white, minimum height=0.8cm},
	% volle Breite
	nbfull/.style={nb, minimum width=\nswidth,
		text width=\dimexpr\nswidth-8pt\relax},
	% eine Ebene eingerueckt (Schleifenrumpf)
	nbi/.style={nb, minimum width=\dimexpr\nswidth-\nsind\relax,
		text width=\dimexpr\nswidth-\nsind-8pt\relax},
	% zwei Ebenen eingerueckt (Schleife in Schleife)
	nbii/.style={nb, minimum width=\dimexpr\nswidth-2\nsind\relax,
		text width=\dimexpr\nswidth-2\nsind-8pt\relax},
	% Verzweigung: halbe Breiten
	nbh/.style={nb, minimum width=\dimexpr\nswidth/2\relax,
		text width=\dimexpr\nswidth/2-8pt\relax},
	nbih/.style={nb, minimum width=\dimexpr(\nswidth-\nsind)/2\relax,
		text width=\dimexpr(\nswidth-\nsind)/2-8pt\relax},
	% Verzweigung halbe Breite, zwei Ebenen eingerueckt
	nbiih/.style={nb, minimum width=\nshalfii,
		text width=\dimexpr\nshalfii-8pt\relax},
	% Kopf einer Verzweigung (das Dreieck), inner sep 0
	nc/.style={draw, line width=0.4pt, anchor=north west, inner sep=0pt,
		minimum height=1.0cm, fill=white},
	ncfull/.style={nc, minimum width=\nswidth},
	nci/.style={nc, minimum width=\dimexpr\nswidth-\nsind\relax},
	ncii/.style={nc, minimum width=\dimexpr\nswidth-2\nsind\relax},
}

% Zeichnet das Dreieck + Beschriftungen in einen bereits gesetzten nc-Knoten.
% #1 Knotenname  #2 Bedingung  #3 ja-Label  #4 nein-Label
\newcommand{\nscond}[4]{%
	\draw[line width=0.4pt] (#1.north west) -- (#1.south);
	\draw[line width=0.4pt] (#1.north east) -- (#1.south);
	\node[font=\scriptsize, anchor=north, inner sep=1pt] at ($(#1.north)+(0,-1.5pt)$) {#2};
	\node[font=\scriptsize, anchor=south west, inner sep=1pt] at ($(#1.south west)+(2.5pt,1.5pt)$) {#3};
	\node[font=\scriptsize, anchor=south east, inner sep=1pt] at ($(#1.south east)+(-2.5pt,1.5pt)$) {#4};
}

% Zeichnet den L-Rahmen einer Schleife (linker + unterer Rand der Einrueckspalte).
% #1 Kopfknoten der Schleife  #2 letzter Rumpfknoten
\newcommand{\nsframe}[2]{%
	\draw[line width=0.4pt] (#1.south west) |- (#2.south west);
}


\title{Probeklausur 2}
\author{Luis Staudt}
\date{}

\begin{document}

% --- Titelblock ---
\begin{center}
	\rule{\textwidth}{0.8pt}\\[0.6em]
	{\LARGE\bfseries Probeklausur 2}\\[0.3em]
	{\large Grundlagen der Programmierung}\\[0.7em]
	\begin{tabular}{r l @{\hspace{1.4cm}} r l}
		\textbf{Studiengänge:}    & PEB \& MVB            & \textbf{Autor:}        & Luis Staudt \\
		\textbf{Semester:}        & Sommersemester 2026   & \textbf{Schwierigkeit:}& einfach \\
		\textbf{Bearbeitungszeit:}& 90 Minuten            & \textbf{Gesamtpunkte:} & 90 \\
	\end{tabular}\\[0.7em]
	\rule{\textwidth}{0.8pt}
\end{center}
\vspace{0.5em}

\noindent\textit{Bearbeiten Sie alle fünf Aufgaben. Achten Sie auf sauberen, kompilierfähigen Java-Quelltext. Hilfsmittel: OpenBook, keine elektronischen Geräte.}

\pagestyle{fancy}
\fancyhf{}
\lhead{\small Probeklausur 2, Grundlagen der Programmierung}
\rhead{\small SoSe 2026 (einfach)}
\cfoot{\thepage}
\renewcommand{\headrulewidth}{0.4pt}

% =============================================
\clearpage
\section*{Aufgabe 1: Struktogramm in Java umsetzen \hfill (24 Punkte)}

Das folgende Struktogramm bildet zunächst die Quersumme einer Zahl $m$ (und zählt deren gerade Ziffern) und gibt danach alle Primzahlen bis zu einer Obergrenze $g$ aus.

Setzen Sie das Struktogramm \textbf{vollständig in ein lauffähiges Java-Programm} um. Verwenden Sie zum Einlesen \texttt{Tastatureingabe.leseInt()}.

\begin{center}
\begin{tikzpicture}
	\node[nbfull] (a1) at (0,0) {lies eine ganze Zahl $m$ ein};
	\node[nbfull] (a2) at (a1.south west) {\texttt{qsumme = 0;}\qquad \texttt{anzG = 0;}};

	% while
	\node[nbfull] (h1) at (a2.south west) {\textbf{solange} $m > 0$};
	\node[nbi]    (w1) at ($(h1.south west)+(\nsind,0)$) {\texttt{rest = m \% 10;}\quad \texttt{qsumme = qsumme + rest;}};
	\node[nci]    (wc) at (w1.south west) {};
	\nscond{wc}{\texttt{rest \% 2 == 0}}{ja}{nein}
	\node[nbih]   (wcl) at (wc.south west) {\texttt{anzG = anzG + 1;}};
	\node[nbih]   (wcr) at ($(wc.south west)+(\nshalfi,0)$) {\textit{(nichts)}};
	\node[nbi]    (w2) at (wcl.south west) {\texttt{m = m / 10;}};
	\nsframe{h1}{w2}

	\node[nbfull] (a3) at ($(w2.south west)+(-\nsind,0)$) {drucke \texttt{qsumme} und \texttt{anzG}};
	\node[nbfull] (a4) at (a3.south west) {lies Obergrenze $g$ ein;\qquad \texttt{anzP = 0;}};

	% for z
	\node[nbfull] (h2) at (a4.south west) {\textbf{für} $z = 2,\,3,\,\ldots,\,g$};
	\node[nbi]    (f1) at ($(h2.south west)+(\nsind,0)$) {\texttt{istPrim = true;}};

	% inner for t
	\node[nbi]    (h3) at (f1.south west) {\textbf{für} $t = 2,\,3,\,\ldots,\,z-1$};
	\node[ncii]   (ic) at ($(h3.south west)+(\nsind,0)$) {};
	\nscond{ic}{\texttt{z \% t == 0}}{ja}{nein}
	\node[nbiih]  (icl) at (ic.south west) {\texttt{istPrim = false;}};
	\node[nbiih]  (icr) at ($(ic.south west)+(\nshalfii,0)$) {\textit{(nichts)}};
	\nsframe{h3}{icl}

	% Verzweigung auf Ebene 1
	\node[nci]    (fc) at ($(icl.south west)+(-\nsind,0)$) {};
	\nscond{fc}{\texttt{istPrim == true}}{ja}{nein}
	\node[nbih]   (fcl) at (fc.south west) {drucke $z$;\quad \texttt{anzP = anzP + 1;}};
	\node[nbih]   (fcr) at ($(fc.south west)+(\nshalfi,0)$) {\textit{(nichts)}};
	\nsframe{h2}{fcl}

	\node[nbfull] (a5) at ($(fcl.south west)+(-\nsind,0)$) {drucke \texttt{anzP}};
\end{tikzpicture}
\end{center}

% =============================================
\clearpage
\section*{Aufgabe 2: Energieformel auswählen \hfill (14 Punkte)}

Je nach physikalischer Situation wird eine andere Energieform berechnet.

\begin{formeln}
\[
	\text{(1) kinetisch:}\quad E = \tfrac{1}{2}\,m\,v^2
	\qquad
	\text{(2) potentiell:}\quad E = m\,g\,h \;\;(g = 9{,}81)
\]
\[
	\text{(3) Federspannung:}\quad E = \tfrac{1}{2}\,D\,s^2
\]
\end{formeln}

Schreiben Sie ein Programm, das den Benutzer die Energieform per Zahl (\texttt{1}, \texttt{2} oder \texttt{3}) wählen lässt und anschließend die benötigten Größen abfragt. Berechnen und drucken Sie die Energie mit der \textbf{richtigen} Formel. Eine Masse $m \le 0$ bzw.\ eine Federkonstante $D \le 0$ ist als ungültige Eingabe abzufangen. Verwenden Sie \texttt{if/else}.

\begin{beispiel}[Beispieldialog]
\begin{lstlisting}[language={}]
Energieform (1=kinetisch, 2=potentiell, 3=Feder): 1
Masse m [kg]: 2
Geschwindigkeit v [m/s]: 3
Energie E = 9.0 J
\end{lstlisting}
\end{beispiel}

% =============================================
\clearpage
\section*{Aufgabe 3: Konstruktor für Polarkoordinaten \hfill (20 Punkte)}

Gegeben ist die aus Vorlesung und Übung bekannte Klasse \texttt{Punkt} (Attribute \texttt{xKoordinate}, \texttt{yKoordinate} als \texttt{double}).

Implementieren Sie für die Klasse \texttt{Punkt} einen \textbf{Konstruktor für Polarkoordinaten}. Erster Parameter ist der Radius~$r$ (Entfernung zum Ursprung), zweiter Parameter der Winkel~$\alpha$ zur x-Achse im mathematisch positiven Sinn, angegeben in \textbf{Grad}.

\begin{formeln}[Umrechnung Polar $\to$ kartesisch ($\alpha$ in Grad)]
\[
	x = r \cdot \cos\!\left(\alpha \cdot \tfrac{\pi}{180}\right)
	\qquad
	y = r \cdot \sin\!\left(\alpha \cdot \tfrac{\pi}{180}\right)
\]
\end{formeln}

\begin{lstlisting}
Punkt p1 = new Punkt(4, 30);     // Radius 4, Winkel 30 Grad
Punkt p2 = new Punkt(7.5, 0);    // auf der x-Achse
Punkt p3 = new Punkt(3, 270);    // auf der negativen y-Achse

double a = p1.berechneAbstandVomNullpunkt();  // muss 4.0 ergeben
\end{lstlisting}

\begin{hinweis}[Merkmale eines Konstruktors]
	Er heißt wie die Klasse, hat \emph{keinen} Rückgabewert (auch kein \texttt{void}) und wird mit \texttt{new} aufgerufen.
\end{hinweis}

% =============================================
\clearpage
\section*{Aufgabe 4: Fachwerk, größte Zug- und Druckkraft \hfill (16 Punkte)}

Gegeben ist die aus Vorlesung und Übung bekannte Datei \texttt{fachwerk.jar}. Für ein gewähltes Fachwerk sollen nach der Berechnung die betragsmäßig \textbf{größte Zugkraft} und die betragsmäßig \textbf{größte Druckkraft} samt zugehöriger Stabnummer ausgegeben werden.

\begin{center}
\begin{tikzpicture}[>=Stealth, scale=1.0,
	knoten/.style={circle, fill=black, inner sep=1.6pt},
	stab/.style={thick, gray!70!black}]
	\coordinate (A) at (0,0); \coordinate (B) at (2,0); \coordinate (C) at (4,0); \coordinate (D) at (6,0);
	\coordinate (E) at (1,1.7); \coordinate (F) at (3,1.7); \coordinate (G) at (5,1.7);
	\draw[stab] (A)--(B); \draw[stab] (B)--(C); \draw[stab] (C)--(D);
	\draw[stab] (A)--(E); \draw[stab] (E)--(B); \draw[stab] (E)--(F);
	\draw[stab] (F)--(B); \draw[stab] (F)--(C); \draw[stab] (F)--(G);
	\draw[stab] (G)--(C); \draw[stab] (G)--(D);
	\foreach \p in {A,B,C,D,E,F,G} \node[knoten] at (\p) {};
	\draw[->, thick] (C)--++(0,-0.9);
	\node[below, font=\footnotesize] at ($(C)+(0,-0.9)$) {Last};
\end{tikzpicture}\\[-0.3em]
{\footnotesize\itshape Schematische Darstellung eines Fachwerks.}
\end{center}

\begin{hinweis}[Benötigte Methoden aus \texttt{fachwerk.jar}]
	\texttt{int gibAnzahlStaebe()}, \texttt{Fachwerkstab gibStab(int i)}, \texttt{double gibStabkraft()} (positiv = Zug, negativ = Druck).
\end{hinweis}

\begin{lstlisting}
public class Fachwerkaufgabe {
    public static void main(String[] args) {
        System.out.println("Welches Fachwerk wollen Sie rechnen (0-4)?");
        int wahl = Tastatureingabe.leseInt();
        Fachwerk fw = new Fachwerk(wahl);
        fw.berechneGroessen();
        // hier erweitern: groesste Zug- und Druckkraft mit Stabnummer suchen
    }
}
\end{lstlisting}

\begin{beispiel}[Mögliche Ausgabe]
\begin{lstlisting}[language={}]
Groesste Zugkraft  (Stab 5): 53.033
Groesste Druckkraft (Stab 3): -62.5
\end{lstlisting}
\end{beispiel}

% =============================================
\clearpage
\section*{Aufgabe 5: Wertetabelle in eine Datei schreiben \hfill (16 Punkte)}

Ein Körper wird aus der Ruhe gleichmäßig beschleunigt. Es gilt
\[ v(t) = a \cdot t \qquad\text{und}\qquad s(t) = \tfrac{1}{2}\,a\,t^2, \qquad a = 2{,}5~\text{m/s}^2. \]

Schreiben Sie ein Programm, das für $t = 0~\text{s}$ bis $t = 12~\text{s}$ in Schritten von $1~\text{s}$ jeweils Geschwindigkeit und Strecke berechnet und als drei Spalten in die Datei \texttt{bewegung.txt} schreibt.

\begin{hinweis}[Hinweis]
	Verwenden Sie \texttt{java.io.FileWriter}: mit \lstinline|new FileWriter("bewegung.txt")| öffnen, je Zeile mit \texttt{write(...)} schreiben (Zeilenumbruch \lstinline|"\n"|), am Ende \texttt{close()}.
	Nutzen Sie \lstinline|"\t"| für die Spaltenabstände.
\end{hinweis}

\begin{beispiel}[Erwarteter Dateiinhalt (Auszug)]
\begin{lstlisting}[language={}]
t [s]   v [m/s]   s [m]
0       0.0       0.0
1       2.5       1.25
2       5.0       5.0
...
12      30.0      180.0
\end{lstlisting}
\end{beispiel}

\end{document}
